\documentclass[11pt,a4paper]{article}

\usepackage[T1]{fontenc}
\usepackage[utf8]{inputenc}
\usepackage[english]{babel}
\usepackage[a4paper,margin=25mm]{geometry}
\usepackage{amsmath,amssymb,amsthm}
\usepackage{graphicx}
\usepackage{booktabs}
\usepackage{multirow}
\usepackage{longtable}
\usepackage{array}
\usepackage{microtype}
\usepackage{xcolor}
\usepackage[hidelinks,colorlinks=false]{hyperref}
\usepackage{enumitem}
\usepackage{caption}

\captionsetup[table]{skip=4pt,labelfont=bf}
\captionsetup[figure]{skip=4pt,labelfont=bf}

\graphicspath{{figs/}}

\newtheorem{lemma}{Lemma}
\newtheorem{hypothesis}{Empirical Hypothesis}
\newtheorem{corollary}{Corollary}

\title{\textbf{Hybrid privacy-aware semantic search: SVD-truncated
document geometry and CKKS-encrypted query reranking under a restricted
threat model}}
\author{Sergey M.\ Kurilenko\\
\small Moscow Institute of Physics and Technology\\
\small \texttt{sergkurilenko@gmail.com}}
\date{\today}

\begin{document}
\maketitle

\begin{abstract}
\noindent
Dense vector embeddings underpin modern semantic search and
Retrieval-Augmented Generation (RAG), yet a growing body of work shows
that an embedding can be inverted back into the text that produced it
with alarming fidelity: once a vector database leaks, the documents
behind it leak with it. The two textbook defences lie at opposite
extremes---encrypting the entire search with fully homomorphic
encryption is cryptographically sound but far too slow for corpora of
millions of documents, whereas injecting privacy-preserving noise
degrades ranking quality long before it provides meaningful protection.

We study a middle path built on an asymmetry between the static document
collection and the dynamic user query. The collection is protected
geometrically: every document vector is compressed onto a
lower-dimensional subspace and then rotated by a secret orthogonal
transformation held only by the data owner. The query is protected
cryptographically: it is processed under homomorphic encryption so that
the server---which honestly follows the protocol but may inspect
everything it stores---never observes the query values or the similarity
scores it computes. The encryption parameters are fixed by a small,
reproducible offline benchmark over the discrete CKKS grid rather than by
manual trial and error.

Beyond the construction itself, we contribute mathematically proven
theoretical results: we derive a tight lower bound on the reconstruction
error of any attacker confined to the protected subspace, making precise
how much information the geometric projection provably removes.
Large-scale experiments on a one-million-document corpus and five
contemporary text encoders corroborate the empirical side---the
protection preserves retrieval quality and, on the strongest encoders,
even slightly improves it by acting as a linear denoiser; query latency
stays comfortably below one second; and an off-the-shelf
text-reconstruction attack confined to the protected space collapses to
the noise floor. We then deliberately step outside that comfortable
regime and measure what happens against stronger adversaries: a
known-plaintext attacker recovers the secret rotation by orthogonal
Procrustes from roughly as many leaked pairs as the retained dimension,
the public product-quantization codes preserve most nearest-neighbour
structure, and a calibrated-noise and random-projection baseline suite
positions the trade-off against natural alternatives; a BEIR check shows
the truncation is an encoder-dependent accuracy cost rather than a free
denoiser on graded-relevance tasks. We therefore state the boundaries
plainly: query confidentiality is cryptographic, but document protection
rests on SVD truncation and a secret rotation that are an empirical
obfuscation layer---not a cryptographic primitive---and we delimit
precisely the threat model under which each claim holds.
\end{abstract}

\textbf{Keywords:}
homomorphic encryption, CKKS, privacy-preserving information retrieval,
semantic search, embedding inversion, non-adaptive attacker, restricted
threat model, SVD truncation, product quantization,
retrieval-augmented generation.

% =====================================================================
\section{Introduction}
% =====================================================================

The privacy of vector representations of text---``embeddings''---has become
a first-class concern in production retrieval systems. State-of-the-art
encoders (Sentence-BERT \cite{ref_sbert}, multilingual-E5 \cite{ref_e5},
GTR \cite{ref_gtr}, BGE-M3 \cite{ref_bgem3}) push the per-document
representation through a
$\sim 10^2$-dimensional bottleneck that nevertheless preserves enough
information to reconstruct the original text with high BLEU.
\textsc{Vec2Text} \cite{ref_vec2text} reports
$\mathrm{BLEU} \approx 97.3\%$ on $32$-token inputs when the attacker has
query access to the encoder; \textsc{GEIA}~\cite{ref_geia} and
\textsc{TEIA}~\cite{ref_teia} show that even snapshot attacks
(an adversary with read access to the vector database) recover personally
identifiable information at high rates. For Retrieval-Augmented Generation
pipelines whose vector stores typically contain customer support tickets,
internal corporate documents and personal data, this rule of thumb---``a
leak of the embedding is a leak of the text''---is now well documented
\cite{ref_zeng_rag}.

The defensive landscape is bracketed by two extremes that, individually,
do not satisfy industrial requirements. Differential privacy
\cite{ref_dwork,ref_dpsgd} adds calibrated noise per-record but, in the
naive coordinate-Gaussian formulation on retrieval-quality embeddings,
destroys ranking quality long before the privacy budget reaches a level
considered useful. Fully homomorphic encryption \cite{ref_ckks,ref_openfhe}
allows arbitrary computation on ciphertexts, but the
ciphertext-ciphertext (ct-ct) regime is too expensive for top-$k$ search
across $10^6$ documents.

\paragraph{Nature and positioning of this work.} This is a
\emph{scientific study}, not a systems-engineering paper. SVD truncation,
random orthogonal rotation, product quantisation and CKKS are
individually well known; we do not claim a new system or a deployable
product. The research questions are: \textbf{(RQ1)} what does projection
provably remove---is there a formal lower bound on what any
projection-restricted decoder can reconstruct? \textbf{(RQ2)} how do SVD
truncation and an unknown secret rotation jointly reshape the
privacy--accuracy--latency trade-off of semantic search at
$10^6$-document scale? \textbf{(RQ3)} does data-driven SVD truncation
behave as a destructive lossy step, or---as we find---as a linear
\emph{denoiser} for retrieval-trained encoders? Accordingly, the secret
rotation is studied as an empirical obfuscation phenomenon, not proposed
as a cryptographic primitive, and the formal lemma is a projection
bound, not an inversion-security theorem.

\paragraph{Object of study and scientific contribution.} The object we
analyse is an \emph{asymmetric} hybrid construction (described below);
the contribution is its theoretical and empirical characterisation under
a credible---but explicitly restricted---threat model, not its
engineering realisation:

\begin{itemize}[leftmargin=18pt]
\item For the \emph{static} side (the database), we apply a deterministic
SVD truncation $V_k$ with $k = d/2$, immediately followed by a Haar-uniform
secret orthogonal rotation $R \in O(k)$. The truncation is chosen so that
the proxy criterion $\sigma_{rec} \geq 0.10$ (Section~\ref{sec:lemma}) is
met uniformly for all tested encoders; the rotation hides the orientation
of the basis from any attacker that does not know $R$. We then publish
a Product-Quantisation artefact (codebook + per-document PQ codes), trained
\emph{in the rotated space}, so that stage-1 candidate filtering can be
performed locally by the client without sending the query to the server.

\item For the \emph{dynamic} side (the query), we apply CKKS encryption.
The protocol is engineered so that the server-side reranking over the
short-list of $K_{\mathrm{cands}}$ candidates only needs the
ciphertext-plaintext (ct-pt) mode---the costly relinearisation step
specific to ct-ct multiplication is avoided---bringing the cryptographic
budget for one query well below one second.

\item The CKKS parameters
$(N_{\mathrm{poly}}, \log_2 Q, d_{\mathrm{proj}}, \Delta)$ are fixed by a
small, reproducible offline micro-benchmark over the discrete parameter grid
(Section~\ref{sec:autotune}): we keep the minimum-latency configuration that
satisfies the conservative security tables and an accuracy tolerance. The
selected configuration ($N_{\mathrm{poly}} = 8192$, $[60, 40, 60]$,
$\Delta = 2^{40}$) yields a $\approx 1.7\times$ speed-up over the TenSEAL
stock setting $[60,40,40,60]$ at the same parameter-table security bound.
This is a deterministic engineering choice over a published grid; we make no
machine-learning claim.
\end{itemize}

\paragraph{Notation.} To avoid the overloading of the symbol $N$, we use
$N_{\mathrm{poly}}$ for the CKKS polynomial-modulus degree,
$N_{\mathrm{docs}}$ for the corpus size, and $K_{\mathrm{cands}}$ for the
stage-1 short-list length throughout.

\paragraph{Theorem and an empirical link.} We prove a tight lower bound on
the $L_2$ reconstruction error of any decoder whose image is constrained
to $\mathrm{span}(V_k)$ (Lemma~\ref{lemma:proj-l2}). The bound makes
$\sigma_{rec}$ a meaningful proxy of the difficulty an inverter faces;
its empirical relation to the BLEU of an off-the-shelf
\textsc{Vec2Text} attack is then formalised as Empirical
Hypothesis~\ref{hyp:bleu} and tested on \textsc{GTR-base} embeddings.

\paragraph{Roadmap.} Section~\ref{sec:related} reviews the related
literature; Section~\ref{sec:threat} introduces the threat model,
separating \emph{document}, \emph{query} and \emph{access-pattern}
privacy, and gives both a coverage table and a claims-vs-evidence table;
Section~\ref{sec:method} describes the asymmetric architecture, including
a cryptographic-reproducibility subsection;
Section~\ref{sec:autotune} states the reproducible CKKS parameter selection;
Section~\ref{sec:lemma} proves the projection-decoder lemma and
delimits its scope; Section~\ref{sec:experiments} presents four
experiments; Section~\ref{sec:discussion} discusses limitations,
required attack experiments and open work.

% =====================================================================
\section{Related Work}\label{sec:related}
% =====================================================================

\paragraph{Embedding inversion.} \textsc{Vec2Text} \cite{ref_vec2text}
trains an encoder--decoder model that iteratively rewrites a candidate
hypothesis until its embedding matches a target vector; on
$32$-token \textsc{MS-MARCO} inputs the attack achieves
$\mathrm{BLEU} \approx 97.3\%$ and $\approx 92\%$ exact recovery.
\textsc{GEIA} \cite{ref_geia} instead trains a generative decoder that
maps a target embedding to a prompt and reconstructs the whole sentence
in a single pass. \textsc{TEIA} \cite{ref_teia} removes the need to query
the victim encoder, training a surrogate inverter that transfers across
embedding models. Adjacent work on
membership-inference \cite{ref_shokri} and on extracting training
texts from large language models \cite{ref_carlini} confirms that the
risk surface of dense retrievers is broader than the embedding-inversion
literature alone suggests.

\paragraph{Modern inversion attacks (2024--2026).} Since \textsc{Vec2Text}
the threat has strengthened along two axes that matter for any geometric
obfuscation. \emph{Few-shot alignment} attacks such as
\textsc{ALGEN}~\cite{ref_algen} learn a linear map between embedding spaces
from $\sim\!10^{3}$ leaked pairs and reach \textsc{Vec2Text}-level recovery
without query access to the victim encoder; \emph{zero-shot} attacks such as
\textsc{ZSInvert}~\cite{ref_zsinvert} invert with no encoder-specific
training; the unsupervised translator \textsc{vec2vec}~\cite{ref_vec2vec}
aligns two embedding spaces with no paired data at all and argues that a
leaked vector store should be treated as plaintext; and
\textsc{Zero2Text}~\cite{ref_zero2text} updates an online token-by-token
regression. These attacks are \emph{alignment-based} or \emph{adaptive} and
exploit the universal geometry of embedding spaces, so a single secret
rotation---which only hides the orientation of a fixed basis---does not by
itself neutralise them. They therefore lie outside the non-adaptive threat
model under which our construction is analysed (Section~\ref{sec:threat});
we quantify the closest members of this family---known-plaintext linear
alignment and an aligned off-the-shelf inverter---directly in
Section~\ref{sec:stronger}, and flag the fully learned variants as the
main residual gap in Section~\ref{sec:limits}.

\paragraph{Differential privacy for embeddings.} Lyu et al.~\cite{ref_lyu}
investigate per-coordinate Gaussian noise on dense representations; the
empirical finding is that retrieval quality degrades sharply long before
$\varepsilon$ reaches a useful range. Random-projection-based privacy
\cite{ref_liu,ref_kenthapadi} predates modern dense retrieval; the
data-dependent SVD projection used here is closer to a noise-removal
technique than to a privacy-preserving DP mechanism, although
$\sigma_{rec}$ provides a useful proxy quantity.

\paragraph{Homomorphic encryption.} The CKKS scheme~\cite{ref_ckks}
introduced approximate-arithmetic FHE; OpenFHE \cite{ref_openfhe} and
TenSEAL provide production-grade implementations.
The bootstrapping cost \cite{ref_ckks_bootstrap} pushes practitioners
to leveled circuits; specialised compilers \cite{ref_eva,ref_chet}
optimise scale and modulus chains.
GPU acceleration of CKKS bootstrap \cite{ref_jung_gpu} and Intel HEXL
\cite{ref_hexl} have made larger leveled circuits feasible.

\paragraph{Hybrid privacy-preserving retrieval.} Tiptoe \cite{ref_tiptoe}
combines clustering with PIR for query privacy. SealPIR
\cite{ref_sealpir}, OnionPIR \cite{ref_onionpir} and SimplePIR
\cite{ref_simplepir} provide single-server PIR with practical query sizes.
The architecture in this paper is complementary: ct-pt CKKS reranking
takes care of the value side of the leakage, while access patterns
remain a separate problem for which PIR/ORAM-like primitives are an
appropriate composition.

% =====================================================================
\section{Threat Model}\label{sec:threat}
% =====================================================================

We adopt the asymmetric threat model
\emph{trusted client / honest-but-curious server / non-adaptive
attacker / unknown rotation}, characterised as follows.

\begin{itemize}[leftmargin=18pt,topsep=4pt,itemsep=2pt]
\item The \emph{client} is the data owner. It runs the encoder $E$,
generates the secret keys $\mu, V_k, R, sk_{\mathrm{CKKS}}$ and never
shares them. The client is trusted in software and key management.
\item The \emph{server} runs the search service over the rotated database
$E_{\mathrm{rot}} \in \mathbb{R}^{N_{\mathrm{docs}} \times k}$ and the
public PQ artefact. It is honest-but-curious: it follows the protocol but
attempts to extract textual information from the data it sees and from the
queries it processes.
\item The \emph{attacker} is non-adaptive: it uses an off-the-shelf
\textsc{Vec2Text} model trained without prior knowledge of $R$. This is
the model under which the construction is \emph{analysed}. We do not stop
there: Section~\ref{sec:stronger} steps outside it and empirically
measures known-plaintext recovery of $R$, public-PQ-code leakage and an
aligned off-the-shelf inverter, so the boundary of the non-adaptive model
is quantified rather than merely asserted. A fully learned decoder
retrained end-to-end against the rotated space, and malicious (rather than
honest-but-curious) server behaviour, remain out of scope.
\end{itemize}

\paragraph{Three distinct privacy notions.} A recurring source of
confusion in hybrid designs is the conflation of \emph{what} is
protected. We therefore separate three independent notions and state,
for each, exactly which component is responsible:

\begin{itemize}[leftmargin=18pt,topsep=4pt,itemsep=2pt]
\item \textbf{Document privacy.} The plaintext rotated database
$E_{\mathrm{rot}}$ \emph{is} visible to the server. It is protected only
by SVD truncation ($\sigma_{rec}$, a proxy quantity) and the secret
rotation $R$ (an empirical obfuscation layer). \emph{This is not a
cryptographic protection}; it fails under known-$R$, known-plaintext,
or an adaptive decoder.
\item \textbf{Query privacy.} CKKS hides the numerical value of the
query vector and of the similarity scores from the server. This part
\emph{is} cryptographic, at the chosen security level.
\item \textbf{Access-pattern privacy.} The candidate IDs reranked per
query, and the public PQ codes, are \emph{not} hidden. CKKS does not
address this channel; it requires composition with a PIR/ORAM-style
primitive.
\end{itemize}

In particular, the public PQ artefact (codebook + per-document codes,
trained in the rotated space) is itself a lossy compressed view of
$E_{\mathrm{rot}}$. Whether neighbourhood, cluster or topic structure
leaks from the public PQ codes alone is measured directly in
Section~\ref{sec:stronger} (Table~\ref{tab:pq-leakage}): it does, and we
treat the PQ artefact as an exposed channel accordingly.

Table~\ref{tab:threat} summarises the coverage and
Table~\ref{tab:claims} maps each claim to the evidence that supports it.
The coverage table is intentionally asymmetric: ct-pt reranking takes
care of the value channel, while the access-pattern channel (which document
IDs are reranked) is left open and would require composition with a
PIR-style primitive, and the public-PQ-artefact and known-plaintext
channels---now quantified in Section~\ref{sec:stronger}---remain exposed
rather than defended.

\begin{table}[t]
\centering
\caption{Threat coverage of the proposed protocol.}\label{tab:threat}
\small
\begin{tabular}{@{}p{0.42\linewidth}p{0.10\linewidth}p{0.42\linewidth}@{}}
\toprule
\textbf{Attack vector} & \textbf{Covered?} & \textbf{Comment} \\ \midrule
Snapshot of $E_{\mathrm{rot}}$ without $R$ (off-the-shelf Vec2Text) &
Partial &
Heuristic; empirically validated against non-adaptive Vec2Text. \\
Network interception, leakage of similarity scores &
Yes &
CKKS at tc128; $sk$ never leaves the client. \\
Known $R$ (known-rotation attacker) &
No &
Reduces to $\sigma_{rec}$ from SVD truncation only. \\
Known-plaintext pairs $(\mathrm{text}_i, E_{\mathrm{rot},i})$ &
No &
$R$ recovered by Procrustes from $\approx k$ pairs
(Sec.~\ref{ssec:procrustes}, Table~\ref{tab:procrustes}). \\
Adaptive decoder retrained against rotated space &
No &
Aligned off-the-shelf inverter measured (Sec.~\ref{ssec:aligned}); a
\emph{learned} decoder is the main open case. \\
Malicious server (protocol deviation, substitution) &
No &
Verifiable computation / TEEs are an orthogonal addition. \\
Access-pattern leakage on candidate IDs &
No &
CKKS hides values, not access patterns; requires PIR/ORAM. \\
Leakage from the public PQ artefact (rotated-space codes) &
No &
Quantified: codes preserve cosine $0.95$ and $67\%$ of top-10 neighbours
(Sec.~\ref{ssec:pqleak}, Table~\ref{tab:pq-leakage}). \\
\bottomrule
\end{tabular}
\end{table}

\begin{table}[t]
\centering
\caption{Claims vs.\ evidence. The table makes explicit which statements
are demonstrated and which are deliberately \emph{not} claimed.}
\label{tab:claims}
\small
\begin{tabular}{@{}p{0.40\linewidth}p{0.54\linewidth}@{}}
\toprule
\textbf{Claim} & \textbf{Evidence / status} \\ \midrule
CKKS hides query values and scores from the server &
Cryptographic construction + parameters (Sec.~\ref{sec:autotune},
\ref{ssec:crypto-detail}); query privacy is cryptographic. \\
ct-pt is faster than ct-ct for the reranking &
Experiment~\ref{ssec:ckks-modes} ($1.44\times$ on the micro-benchmark,
$1.7\times$ at the auto-tuned configuration). \\
Secret rotation helps against \emph{off-the-shelf}, non-adaptive
Vec2Text & Experiment~\ref{ssec:vec2text}; effect holds only for the
tested weak-attacker configuration. \\
The protective wrapper preserves ranking in $\mathrm{span}(V_k)$ &
Experiment~\ref{ssec:integral} (within the 5-seed CI). \\
End-to-end query latency $<1$~s at $10^6$ docs &
Experiment~\ref{ssec:latency} (loopback PoC). \\
The system is secure against an \emph{adaptive} inversion attacker &
\textbf{Not claimed}; aligned off-the-shelf inverter tested
(Sec.~\ref{ssec:aligned}), learned decoder open (Sec.~\ref{sec:limits}). \\
The rotation is a cryptographic guarantee &
\textbf{Not claimed}; it is an empirical obfuscation layer only. \\
Document privacy on the server is cryptographic &
\textbf{Not claimed}; $E_{\mathrm{rot}}$ is plaintext on the server. \\
No information leaks from the public PQ codes &
\textbf{Not claimed}; codes leak neighbour structure
(Sec.~\ref{ssec:pqleak}, Table~\ref{tab:pq-leakage}). \\
\bottomrule
\end{tabular}
\end{table}

% =====================================================================
\section{Method}\label{sec:method}
% =====================================================================

The asymmetric defence is illustrated in Figure~\ref{fig:architecture}.
The architecture is the union of an offline data-preparation phase
(Section~\ref{ssec:offline}) and the online query protocol
(Section~\ref{ssec:online}).

\begin{figure}[tb]
\centering
\includegraphics[width=\linewidth]{architecture_en.pdf}
\caption{Online query flow of the construction under study.
Steps 1--7 realise the two-stage protocol: client transformation
$T(\cdot)$ and local PQ filtering produce a short-list of
$K_{\mathrm{cands}}$ candidate IDs; step~4 encrypts the rotated query
under CKKS; step~5
runs ct-pt reranking on the server; steps~6--7 decrypt the scores and
sort. The secret keys ($\mu$, $V_k$, $R$, $sk_{\mathrm{CKKS}}$) and the
rotated database $E_{\mathrm{rot}}$ are produced offline by the data
owner (Section~\ref{ssec:offline}).}
\label{fig:architecture}
\end{figure}

% ----------------------------------------------------------------
\subsection{Offline phase}\label{ssec:offline}
% ----------------------------------------------------------------

\paragraph{Encoder and centring.} The data owner runs the encoder $E$ on a
corpus $D = \{T_1, \dots, T_{N_{\mathrm{docs}}}\}$ to obtain
$X \in \mathbb{R}^{N_{\mathrm{docs}} \times d}$ with $L_2$-normalised rows.
The global centroid $\mu \in \mathbb{R}^d$ is computed and the corpus is
centred: $X_{\mathrm{c}} = X - \mu$.

\paragraph{SVD truncation.} A randomised SVD~\cite{ref_halko}
$X_{\mathrm{c}} \approx U_k \Sigma_k V_k^\top$ produces the orthonormal
matrix $V_k \in \mathbb{R}^{d \times k}$ that spans the dominant
$k$-dimensional subspace. We use $k = d/2$ in the canonical operating
point: this is the smallest $k$ at which the relative reconstruction
error
\[
\sigma_{rec}^{2}(E; V_k) =
\frac{\|E - \pi_k(E)\|_F^2}{\|E\|_F^2} = 1 - \eta_k
\]
exceeds the proxy threshold $\sigma_{rec} \geq 0.10$ uniformly across
all five tested encoders (Section~\ref{sec:experiments}).

\paragraph{Secret rotation.} A Haar-uniform random orthogonal matrix
$R \in O(k)$ is generated by QR-decomposing a Gaussian matrix and
correcting the diagonal of $R$ to obtain a uniform distribution on the
group. The protected document representation is
\[
v'_i \;=\; T(v_i) \;=\; R\, V_k^\top (v_i - \mu) \in \mathbb{R}^k .
\]

\paragraph{Public PQ artefact.} A faiss \texttt{IndexPQ} index is trained
\emph{in the rotated space} $E_{\mathrm{rot}}$ with $M = k/4$ subquantisers
and $8$ bits per code. Both the codebook and the per-document codes are
public; the client downloads them once at on-boarding. Training the PQ
artefact in the rotated space removes one trivial cross-space pairing
(an attacker cannot compose pairs
$(\hat E_{\mathrm{proj}}, E_{\mathrm{rot}})$ from a non-rotated PQ index
to estimate $R$). It does \emph{not}, however, make the PQ codes safe:
the public codes are a lossy quantised image of $E_{\mathrm{rot}}$ and
may themselves leak neighbourhood, cluster or topic structure. We treat
the PQ artefact as part of the attack surface and measure this leakage
directly in Section~\ref{sec:stronger} (Table~\ref{tab:pq-leakage}): the
public codes preserve mean cosine $0.95$ and most top-10 neighbours, so
they must be treated as an exposed channel.

\paragraph{Client-side state.} The client retains $\mu$, $V_k$, $R$ and
the CKKS secret key $sk_{\mathrm{CKKS}}$. The server stores
$E_{\mathrm{rot}} \in \mathbb{R}^{N_{\mathrm{docs}} \times k}$ in plaintext
and the public PQ artefact.

% ----------------------------------------------------------------
\subsection{Online query protocol}\label{ssec:online}
% ----------------------------------------------------------------

The seven-step protocol of Figure~\ref{fig:architecture} processes a query
in time pipeline:

\begin{enumerate}[leftmargin=22pt,topsep=2pt,itemsep=2pt]
\item \emph{Encoder.} The client computes
$v_q = E(q_{\mathrm{text}}) \in \mathbb{R}^d$.
\item \emph{Transformation.} The rotated query is
$q' = R\, V_k^\top (v_q - \mu) \in \mathbb{R}^k$.
\item \emph{Local PQ search.} Using the public PQ artefact and the
rotated query $\tilde q = q'$, the client performs an asymmetric
PQ-distance search and returns the top-$K_{\mathrm{cands}}$ candidate
IDs $I_{\mathrm{cand}}$. The candidate set is small enough
($K_{\mathrm{cands}} = 40$ in our experiments) that a CKKS reranking
fits within the latency budget.
\item \emph{CKKS encryption.} The client encrypts $q'$:
$\mathrm{ct}_q = \mathrm{Enc}_{pk}(q')$.
\item \emph{Server-side reranking.} For each $j \in I_{\mathrm{cand}}$
the server computes
$\mathrm{ct}_{\mathrm{score}}^{(j)} = \mathrm{ct}_q \odot v'_j$
in the ct-pt mode. The result is a ciphertext encoding the inner
product $\langle q', v'_j \rangle$ in the rotated projected space.
\item \emph{Decryption.} The client decrypts
$s_j = \mathrm{Dec}_{sk}(\mathrm{ct}_{\mathrm{score}}^{(j)})$.
\item \emph{Sort.} The decrypted scores are sorted and the top-$K$
documents returned.
\end{enumerate}

\paragraph{Why ct-pt is sufficient.} In step~5 the server multiplies a
ciphertext by a plaintext vector. Unlike ciphertext-ciphertext
multiplication, this operation produces a ciphertext that stays in the
two-component form, no relinearisation is required, and the modulus chain
is consumed by exactly one rescale. The sum of slots into a single
scalar $\langle q', v'_j \rangle$ requires $\log_2 k$ rotations using
Galois keys, but no further multiplications. The result is a sub-second
latency at $N_{\mathrm{docs}} = 10^6$.

\paragraph{Choice of $K_{\mathrm{cands}}$.} Since each of the
$K_{\mathrm{cands}}$ candidates is processed by one ct-pt operation, the
server-side latency grows linearly in $K_{\mathrm{cands}}$. Conversely, a
small $K_{\mathrm{cands}}$ trades against PQ-recall: the relevant
document must be in the short-list for the CKKS reranker to score it.
We pick $K_{\mathrm{cands}}$ as the smallest value for which
PQ-recall@$K_{\mathrm{cands}}$ matches the SVD-projected exact baseline
within the 5-seed CI; in our integral experiment this is
$K_{\mathrm{cands}} = 40$, but the parameter is a deployment knob and
should be retuned for new encoders or larger corpora.

% ----------------------------------------------------------------
\subsection{Cryptographic reproducibility}\label{ssec:crypto-detail}
% ----------------------------------------------------------------

For the selected configuration
($N_{\mathrm{poly}} = 8192$, coefficient-modulus chain $[60,40,60]$ bits,
$\log_2 Q = 160$, scale $\Delta = 2^{40}$, TenSEAL/Microsoft SEAL
backend) we report the parameters needed to reproduce the cryptographic
budget rather than only the latency:

\begin{itemize}[leftmargin=18pt,topsep=2pt,itemsep=2pt]
\item \emph{Packing.} A single rotated query $q' \in \mathbb{R}^{k}$ is
packed into one ciphertext using $k \leq N_{\mathrm{poly}}/2 = 4096$
slots; for the integral encoders $k \in \{192,384,512\}$, so one
ciphertext per query suffices and no cross-ciphertext aggregation is
needed.
\item \emph{Operations per score.} Each candidate score is one ct-pt
multiply, one rescale (one level consumed) and a slot-sum implemented as
$\lceil \log_2 k \rceil$ ciphertext rotations using Galois keys
($8$--$9$ rotations for the tested $k$); no relinearisation and no
ciphertext-ciphertext multiply occur.
\item \emph{Level/noise budget.} The chain $[60,40,60]$ provides one
multiplicative level beyond the input; after the single rescale the
remaining modulus ($60$ bits) keeps the additive noise far below the
$\Delta = 2^{40}$ scale, so decrypted scores match the plaintext inner
product to a relative error $<10^{-3}$ (Pearson $>0.9999$ vs.\ exact,
Experiment~\ref{ssec:integral}).
\item \emph{Object sizes.} At $N_{\mathrm{poly}}=8192$, $\log_2 Q=160$:
one fresh ciphertext is $\approx 0.21$~MB; the Galois-key set for the
required power-of-two rotations is $\approx 6$--$8$~MB and the
relinearisation key is \emph{not} generated (ct-pt only); the public key
is $\approx 0.4$~MB. Per query the client uploads one ciphertext
($\approx 0.21$~MB) and downloads $K_{\mathrm{cands}}$ score ciphertexts
($\approx 40 \times 0.21 \approx 8.4$~MB before base64; the measured
on-wire JSON+base64 sizes are reported in
Section~\ref{ssec:latency}).
\item \emph{Security level.} We do not assert a generic ``standard''
label; we report $\log_2 Q = 160 < 218$ for $N_{\mathrm{poly}}=8192$,
which is within the conservative ternary-secret \texttt{tc128} bound of
the HomomorphicEncryption.org tables \cite{ref_he_std}, cross-checked
with the lattice-estimator methodology of Albrecht et al.\
\cite{ref_lwe_estimator}. The standard documents are security
\emph{guidelines}; the concrete estimator/table used is stated so the
claim can be re-derived.
\end{itemize}

% =====================================================================
\section{CKKS Parameter Selection}\label{sec:autotune}
% =====================================================================

The CKKS parameter space is large and discrete:
$N_{\mathrm{poly}} \in \{2^{12}, \dots, 2^{14}\}$, multiple
coefficient-modulus chains of total bit budget bounded by the
conservative security tables of HomomorphicEncryption.org
\cite{ref_he_std} (cross-checked with the lattice estimator
\cite{ref_lwe_estimator}), several scale exponents
$\log_2 \Delta \in \{15, 20, 25, 30, 40\}$ and several projection
dimensions $d_{\mathrm{proj}}$. Analytic complexity is proportional to
$N_{\mathrm{poly}} \log N_{\mathrm{poly}}$, but cache effects, SIMD
vectorisation and slot-packing ratios
($d_{\mathrm{proj}} / (N_{\mathrm{poly}}/2)$) cause step-like deviations
from the analytic prediction, so we fix the configuration by a small
\emph{offline micro-benchmark} rather than analytically.

\paragraph{Selection procedure (reproducible).} We enumerate the discrete
grid; discard configurations that violate the security tables or the
correctness (noise-budget) constraint; time one batched ct-pt operation with
rescale for each survivor; and keep the minimum-latency configuration whose
retrieval accuracy stays within a tolerance $\tau$ of the SVD-projected
baseline, $|\Delta\mathrm{Acc@1}| \leq \tau$. This is a deterministic
engineering search over a published grid, reported only so the chosen budget
can be reproduced; we make \emph{no} machine-learning contribution. With
$\tau = 1$~p.p.\ the procedure selects $N_{\mathrm{poly}} = 8192$,
$[60, 40, 60]$, $\Delta = 2^{40}$ ($\log_2 Q = 160$ bits, within the
conservative $218$-bit \texttt{tc128} bound for $N_{\mathrm{poly}} = 8192$
\cite{ref_he_std,ref_lwe_estimator}), which is $\approx 1.7\times$ faster than
the TenSEAL stock setting $[60, 40, 40, 60]$ ($\log_2 Q = 200$~bits) at the
same parameter-table security bound (Table~\ref{tab:ckks-configs}). The full
cryptographic budget of this configuration is given in
Section~\ref{ssec:crypto-detail}.

\begin{table}[tb]
\centering
\caption{CKKS configurations on the test workload (one batched ct-pt
reranking operation). All rows sit within the conservative \texttt{tc128}
($\geq 128$-bit) bound of the HomomorphicEncryption.org tables for the stated
$N_{\mathrm{poly}}$.}
\label{tab:ckks-configs}
\small
\begin{tabular}{@{}lrrrrr@{}}
\toprule
Configuration & $N_{\mathrm{poly}}$ & $\log_2 Q$ (bits) & Security & Time (ms) & Speed-up \\ \midrule
Default (TenSEAL stock, $[60, 40, 40, 60]$) & 8192 & 200 & tc128 ($\geq 128$) & 419.4 & $1.0\times$ \\
Min-latency (Acc@10), $[40, 20, 40]$, $\Delta = 2^{20}$ & 4096 & 100 & tc128 & 173.7 & $2.4\times$ \\
\textbf{Selected (Acc@1), $[60, 40, 60]$, $\Delta = 2^{40}$} & 8192 & 160 & tc128 & 245.2 & $1.7\times$ \\ \bottomrule
\end{tabular}
\end{table}

% =====================================================================
\section{Theoretical Analysis: Decoder Lower Bound}\label{sec:lemma}
% =====================================================================

Let $\mathbf{x} \in \mathbb{R}^d$ denote a centred embedding,
$V_k \in \mathbb{R}^{d \times k}$ the orthonormal matrix from the SVD
truncation, and $\pi_k(\mathbf{x}) = V_k V_k^\top \mathbf{x}$ the
orthogonal projection onto $\mathrm{span}(V_k)$.
Define
$\sigma_{rec}(\mathbf{x}; V_k) =
\| \mathbf{x} - \pi_k(\mathbf{x}) \|_2 / \|\mathbf{x}\|_2$
as the per-vector relative reconstruction error and
$\mathbf{x}_\perp = \mathbf{x} - \pi_k(\mathbf{x})$.

\begin{lemma}\label{lemma:proj-l2}
Let $f : \mathbb{R}^d \to \mathbb{R}^d$ be any decoder whose image is
constrained by $f(\mathbf{y}) \in \mathrm{span}(V_k)$ for every
$\mathbf{y} \in \mathbb{R}^d$. Then for every
$\mathbf{x} \in \mathbb{R}^d$
\[
\| \mathbf{x} - f(\pi_k(\mathbf{x})) \|_2^{2}
\;\geq\; \| \mathbf{x}_\perp \|_2^{2},
\]
with equality at $f(\mathbf{y}) = \mathbf{y}$. In relative form,
$\| \mathbf{x} - f(\pi_k(\mathbf{x})) \|_2 / \|\mathbf{x}\|_2
\geq \sigma_{rec}(\mathbf{x}; V_k)$.
\end{lemma}

\begin{proof}
Decompose $\mathbf{x} = \pi_k(\mathbf{x}) + \mathbf{x}_\perp$ with
$\pi_k(\mathbf{x}) \in \mathrm{span}(V_k)$ and
$\mathbf{x}_\perp \in \mathrm{span}(V_k)^\perp$. By the image
restriction $f(\pi_k(\mathbf{x})) \in \mathrm{span}(V_k)$, hence
$\mathbf{x} - f(\pi_k(\mathbf{x})) =
\bigl( \pi_k(\mathbf{x}) - f(\pi_k(\mathbf{x})) \bigr) + \mathbf{x}_\perp$
is the sum of two orthogonal vectors. Pythagoras gives
$\| \mathbf{x} - f(\pi_k(\mathbf{x})) \|_2^{2}
= \| \pi_k(\mathbf{x}) - f(\pi_k(\mathbf{x})) \|_2^{2}
+ \| \mathbf{x}_\perp \|_2^{2}
\geq \| \mathbf{x}_\perp \|_2^{2}$.
\end{proof}

\paragraph{Scope: a projection bound, not an inversion-security
theorem.} Lemma~\ref{lemma:proj-l2} bounds the $L_2$ error of decoders
\emph{whose image is contained in} $\mathrm{span}(V_k)$. It is a
statement about information lost to the projection, not about the
security of text inversion. A realistic attacker is not so constrained:
it may output arbitrary text, whose re-embedding generally has a
non-zero $\mathbf{x}_\perp$ component, and it may exploit corpus priors
or memorisation to partially recover the discarded component. The lemma
therefore does \emph{not} lower-bound the achievable inversion BLEU,
token overlap or PII recovery. We use it only to make the
\emph{proxy} criterion $\sigma_{rec} \geq 0.10$ precise within the
projection-restricted class, and we stress that the threshold $0.10$ is
an \emph{engineering proxy}, calibrated on one encoder (GTR-base) against
one off-the-shelf attacker, \emph{not} a security threshold. A
per-encoder ablation of $\sigma_{rec}$ against BLEU, token overlap and
typed-PII recovery (names, addresses, e-mail, phone, medical terms) is
listed as a remaining open experiment (Section~\ref{sec:limits}).

\begin{hypothesis}\label{hyp:bleu}
For an off-the-shelf inversion attack \textsc{Vec2Text}~\cite{ref_vec2text},
the expected BLEU of recovering the original text from
$\pi_k(\mathbf{x})$ is monotone in
$\eta_k = 1 - \sigma_{rec}^{2}(E; V_k)$:
\[
\mathbb{E}\bigl[ \mathrm{BLEU}(f(\pi_k(\mathbf{x})), T(\mathbf{x})) \bigr]
\;\approx\; \mathrm{BLEU}_0 + \gamma_f \, \eta_k,
\]
where $\mathrm{BLEU}_0$ is the BLEU of a ``random semantically close''
text and $\gamma_f$ is a decoder-specific constant.
\end{hypothesis}

Hypothesis~\ref{hyp:bleu} is validated numerically in
Section~\ref{ssec:vec2text}; we deliberately separate the formal lemma
from the empirical observation.

% =====================================================================
\section{Experiments}\label{sec:experiments}
% =====================================================================

We report four experiments:
(1) CKKS modes (ct-pt vs ct-ct);
(2) off-the-shelf Vec2Text attack against SVD-truncated (and possibly
rotated) embeddings;
(3) the integral multi-encoder experiment over a $10^6$-document
Russian-Wikipedia corpus;
(4) end-to-end latency on a client-server PoC.

\paragraph{Hardware and software.} All experiments run on a single
workstation: Intel Core i5-14400F, 32~GB DDR5-4800, NVIDIA RTX~5060
(8~GB VRAM), Windows 11, Python 3.11. To make the runs reproducible we
pin exact versions and record the commit/build hashes of the
security-relevant libraries in the supplementary
\texttt{environment.lock}: PyTorch 2.10 + CUDA 12.8,
TenSEAL 0.3.16 (Microsoft SEAL 4.1 backend),
faiss-cpu 1.13.0, transformers 4.57.0, scikit-learn 1.5.x,
sacrebleu for BLEU \cite{ref_bleu}. The pinned hashes matter
because CKKS noise behaviour and faiss PQ training depend on the
backend build.

% ----------------------------------------------------------------
\subsection{Experiment 1: ciphertext-plaintext speed-up}\label{ssec:ckks-modes}
% ----------------------------------------------------------------

We benchmark a batched dot-product of an encrypted $192$-dimensional
query with $N = 10\,000$ plaintext documents.
The CKKS parameters are $N = 8192$,
$\mathrm{coeff\_mod\_bit\_sizes} = [60, 40, 40, 60]$,
$\Delta = 2^{40}$ (security level $\geq 128$ bits).
Two regimes are compared: ct-ct (both query and database encrypted, with
relinearisation per multiplication) and ct-pt (only the query is encrypted,
no relinearisation). Each is run for $50$ independent trials and
the mean is reported with a $95\%$ Gaussian confidence interval.

\begin{table}[tb]
\centering
\caption{Ciphertext-plaintext vs.\ ciphertext-ciphertext for the dot
product against $N = 10\,000$ documents (mean of $50$ trials, $95\%$~CI).}
\label{tab:ckks-modes}
\small
\begin{tabular}{@{}llrrrr@{}}
\toprule
Mode & Operation & Relin.\ & Latency $T_{\mathrm{query}}$ & QPS & Speed-up \\ \midrule
ct-ct (baseline) & $\mathrm{Enc}(v_1)\cdot \mathrm{Enc}(v_2)$ & yes (heavy) & $19.15 \pm 0.15$~s & $0.052$ & $1.00\times$ \\
\textbf{ct-pt (proposed)} & $\mathrm{Enc}(v_1)\cdot \mathrm{Plain}(v_2)$ & no & $\mathbf{13.28 \pm 0.10}$~s & $0.075$ & $\mathbf{1.44\times}$ \\
\bottomrule
\end{tabular}
\end{table}

\begin{figure}[tb]
\centering
\includegraphics[width=0.6\linewidth]{fig_exp2_ctpt_vs_ctct.png}
\caption{ct-pt vs.\ ct-ct: distribution of per-batch latencies over
$50$ trials. Removing relinearisation moves the median by
$\sim 6$~s and tightens the upper tail.}
\label{fig:exp2}
\end{figure}

% ----------------------------------------------------------------
\subsection{Experiment 2: Vec2Text attack}\label{ssec:vec2text}
% ----------------------------------------------------------------

We instantiate \textsc{Vec2Text} on \textsc{GTR-base} embeddings
($d = 768$). The attack is applied to a synthetic 100-document corpus
containing PII (names, addresses, card numbers, medical phrasing).
Three threat models are compared at five values of $k/d$:
(a) no rotation; (b) rotation with the matrix $R$ \emph{known} to the
attacker; (c) rotation with $R$ \emph{unknown}.

The implementation budget of the attacker is fixed (greedy decoding,
$10$ corrector iterations, no beam search, $8$~GB GPU); under this
budget the unprotected baseline reaches
$\mathrm{BLEU}_0 = 0.156$ instead of the $0.973$ reported by
Morris et al.~\cite{ref_vec2text}, because the postulated PII corpus,
greedy decoding and beam-free regime are intentionally weaker than
\cite{ref_vec2text}'s strongest configuration; the experiment measures
the relative effect of SVD truncation and secret rotation in a fixed
weak-attacker configuration.

\paragraph{Strength of the attacker: an explicit caveat.} This is the
single most important limitation of \emph{this particular experiment}: it
establishes the effect of $R$ \emph{against an artificially weak inverter
only}. A stronger attacker---beam search with a large width, many more
corrector iterations, an adaptive decoder fine-tuned on
$(\text{text}, E_{\mathrm{rot}})$ pairs, or a modern universal/zero-shot
inversion model---is not evaluated \emph{here}. We therefore do not rest
the security narrative on this table. Section~\ref{sec:stronger} measures
the two decisive stronger channels: a known-plaintext recovery of $R$ by
orthogonal Procrustes~\cite{ref_schonemann} (Table~\ref{tab:procrustes}),
which succeeds from $\approx k$ leaked pairs, and an aligned off-the-shelf
inverter with a larger budget (Table~\ref{tab:adaptive-inversion}). Every
sentence about the rotation in isolation should accordingly be read as
``against the tested off-the-shelf, non-adaptive Vec2Text configuration'',
with the boundary now quantified in Section~\ref{sec:stronger}.

\begin{table}[tb]
\centering
\caption{Off-the-shelf Vec2Text on GTR-base embeddings under the three
threat models. BLEU and Token Overlap with 95\% bootstrap CIs in
brackets.}
\label{tab:vec2text}
\small
\begin{tabular}{@{}rllll@{}}
\toprule
$k/d$ & Rotation & Knowledge of $R$ & BLEU & Token Overlap \\ \midrule
$1.00$ & off & --- & $0.156\,[0.123;\,0.186]$ & $0.385\,[0.351;\,0.417]$ \\
$1.00$ & on  & known & $0.156\,[0.127;\,0.187]$ & $0.385\,[0.351;\,0.416]$ \\
$1.00$ & on  & \textbf{unknown} & $\mathbf{0.007\,[0.006;\,0.008]}$ & $\mathbf{0.053\,[0.046;\,0.060]}$ \\
$0.75$ & off & --- & $0.077\,[0.063;\,0.093]$ & $0.325\,[0.304;\,0.345]$ \\
$0.75$ & on  & unknown & $0.007\,[0.007;\,0.008]$ & $0.061\,[0.054;\,0.067]$ \\
$0.50$ & off & --- & $0.057\,[0.046;\,0.069]$ & $0.285\,[0.270;\,0.300]$ \\
$0.50$ & on  & unknown & $0.007\,[0.007;\,0.008]$ & $0.050\,[0.045;\,0.055]$ \\
$0.25$ & off & --- & $0.024\,[0.020;\,0.028]$ & $0.196\,[0.181;\,0.208]$ \\
$0.25$ & on  & unknown & $0.007\,[0.006;\,0.007]$ & $0.045\,[0.040;\,0.051]$ \\
$0.10$ & off & --- & $0.010\,[0.009;\,0.011]$ & $0.113\,[0.104;\,0.123]$ \\
$0.10$ & on  & unknown & $0.007\,[0.007;\,0.008]$ & $0.053\,[0.047;\,0.059]$ \\
\bottomrule
\end{tabular}
\end{table}

\begin{figure}[tb]
\centering
\includegraphics[width=0.49\linewidth]{fig_exp3_bleu_vs_k.png}\hfill
\includegraphics[width=0.49\linewidth]{fig_exp3_bleu_kd1.png}
\caption{Vec2Text BLEU as a function of $k/d$ under three threat models
(left); zoom on $k/d = 1.0$ comparing no-rotation, known-$R$ and
unknown-$R$ (right). Lines for ``no rotation'' and ``known-$R$'' overlap;
the ``unknown-$R$'' line stays at the noise level.}
\label{fig:exp3}
\end{figure}

\paragraph{Findings.}
\begin{itemize}[leftmargin=18pt,topsep=2pt]
\item A known $R$ does not provide protection beyond SVD truncation:
the columns ``no rotation'' and ``known-$R$'' are statistically
indistinguishable.
\item Under the tested off-the-shelf configuration, an unknown $R$
reduces BLEU to the observed noise floor ($\sim 0.007$) for all $k/d$
values, including $k/d = 1.0$ (no truncation at all). The ratio
$\mathrm{BLEU}_{\mathrm{off}} / \mathrm{BLEU}_{\mathrm{unknown}}$
is $22\times$ at $k/d = 1.0$, $11\times$ at $0.75$, $8\times$ at $0.50$.
This is an effect against a \emph{non-adaptive} attacker only and is not
evidence of robustness against an adaptive or known-plaintext adversary.
\item Empirical Hypothesis~\ref{hyp:bleu} is consistent with the
no-rotation column: a linear regression of BLEU on
$\eta_k = 1 - \sigma_{rec}^2$ gives $R^{2} \approx 0.93$.
This is consistent with, but does not formally derive,
the BLEU~$\to$~$\sigma_{rec}$ link.
\item In the operating point of the integral experiment ($k/d = 0.5$,
rotation enabled, unknown $R$) the expected BLEU is at the noise level.
\end{itemize}

% ----------------------------------------------------------------
\subsection{Experiment 3: integral multi-encoder evaluation}\label{ssec:integral}
% ----------------------------------------------------------------

We run the full pipeline on the canonical operating point $k = d/2$
across five encoders:

\begin{itemize}[leftmargin=18pt,topsep=2pt,itemsep=2pt]
\item \texttt{intfloat/multilingual-e5-small} ($d = 384$),
\item \texttt{intfloat/multilingual-e5-base} ($d = 768$),
\item \texttt{sentence-transformers/paraphrase-multilingual-mpnet-base-v2} ($d = 768$),
\item \texttt{intfloat/multilingual-e5-large} ($d = 1024$),
\item \texttt{BAAI/bge-m3} ($d = 1024$).
\end{itemize}

The corpus is one million Russian-Wikipedia paragraphs; queries are
$500$ self-retrieval requests sampled at seed $42$;
the rotation is averaged over $5$ seeds
$\{11, 23, 47, 31, 53\}$. PQ uses $M = k/4$ subquantisers with $8$ bits
each, yielding artefacts of $48$/$96$/$128$~MB. The CKKS configuration
is the auto-selected one
($N_{\mathrm{poly}} = 8192$, $[60, 40, 60]$, $\Delta = 2^{40}$).

\paragraph{On the use of self-retrieval.} Self-retrieval (the query is
the first sentence of the target paragraph) is a \emph{controlled
geometry probe}, not a production IR benchmark: it isolates whether the
defensive transform $T(\cdot)$ preserves the ranking induced by the
encoder, with a deterministic, leakage-free ground truth at
$10^6$-scale. It is sufficient for the paper's actual claim---that the
PQ+CKKS wrapper is metric-preserving in $\mathrm{span}(V_k)$---but it
does \emph{not} establish production retrieval quality. Re-running the
integral experiment on standard IR benchmarks (BEIR, MIRACL, MS~MARCO)
is listed as a generality experiment in Section~\ref{sec:limits}; the
present numbers should be read as a transform-fidelity measurement,
not as a retrieval-quality benchmark.

\begin{table}[!t]
\centering
\caption{Integral experiment at $k = d/2$ on a $10^6$-document
corpus across five encoders. Reported are the raw baseline (no
defence), the SVD-projected baseline (FAISS-IP exact in
$\mathrm{span}(V_k)$, no rotation, no CKKS), the proposed pipeline
(mean $\pm 95\%$ t-CI half-width over 5 rotation seeds), and the
delta of Acc@1 from the raw baseline.}
\label{tab:integral}
\small
\setlength{\tabcolsep}{4pt}
\begin{tabular}{@{}llrrrrr@{}}
\toprule
Encoder & Metric & Baseline (raw $d$) & Baseline\_proj (SVD $k$) & Proposed & $\Delta$ from raw & Server $p_{95}$ (ms) \\ \midrule
\multicolumn{7}{l}{\emph{multilingual-e5-small ($d=384$, $k=192$, $M=48$, $\sigma_{rec}=0.239$)}} \\
& Acc@1   & $0.782$ & $0.702$ & $0.701 \pm 0.001$ & $-0.081$ & \multirow{4}{*}{$314$} \\
& Acc@10  & $0.890$ & $0.834$ & $0.831 \pm 0.002$ & $-0.059$ & \\
& MRR     & $0.817$ & $0.747$ & $0.746 \pm 0.001$ & $-0.071$ & \\
& NDCG@10 & $0.834$ & $0.768$ & $0.767 \pm 0.001$ & $-0.067$ & \\ \midrule
\multicolumn{7}{l}{\emph{multilingual-e5-base ($d=768$, $k=384$, $M=96$, $\sigma_{rec}=0.129$)}} \\
& Acc@1   & $0.830$ & $0.834$ & $0.833 \pm 0.001$ & $+0.003$ & \multirow{4}{*}{$356$} \\
& Acc@10  & $0.934$ & $0.924$ & $0.924 \pm 0.001$ & $-0.010$ & \\
& MRR     & $0.862$ & $0.865$ & $0.864 \pm 0.001$ & $+0.002$ & \\
& NDCG@10 & $0.880$ & $0.879$ & $0.879 \pm 0.000$ & $-0.001$ & \\ \midrule
\multicolumn{7}{l}{\emph{paraphrase-multilingual-mpnet-base-v2 ($d=768$, $k=384$, $M=96$, $\sigma_{rec}=0.107$)}} \\
& Acc@1   & $0.716$ & $0.650$ & $0.649 \pm 0.001$ & $-0.067$ & \multirow{4}{*}{$351$} \\
& Acc@10  & $0.842$ & $0.796$ & $0.795 \pm 0.001$ & $-0.047$ & \\
& MRR     & $0.762$ & $0.693$ & $0.692 \pm 0.001$ & $-0.070$ & \\
& NDCG@10 & $0.782$ & $0.717$ & $0.717 \pm 0.000$ & $-0.066$ & \\ \midrule
\multicolumn{7}{l}{\emph{multilingual-e5-large ($d=1024$, $k=512$, $M=128$, $\sigma_{rec}=0.101$)}} \\
& Acc@1   & $0.796$ & $0.814$ & $0.814 \pm 0.000$ & $+0.018$ & \multirow{4}{*}{$219$} \\
& Acc@10  & $0.880$ & $0.922$ & $0.920 \pm 0.000$ & $+0.040$ & \\
& MRR     & $0.824$ & $0.848$ & $0.848 \pm 0.000$ & $+0.024$ & \\
& NDCG@10 & $0.838$ & $0.866$ & $0.865 \pm 0.000$ & $+0.027$ & \\ \midrule
\multicolumn{7}{l}{\emph{BAAI/bge-m3 ($d=1024$, $k=512$, $M=128$, $\sigma_{rec}=0.128$)}} \\
& Acc@1   & $0.818$ & $0.832$ & $0.832 \pm 0.000$ & $+0.014$ & \multirow{4}{*}{$221$} \\
& Acc@10  & $0.908$ & $0.926$ & $0.924 \pm 0.000$ & $+0.016$ & \\
& MRR     & $0.848$ & $0.863$ & $0.863 \pm 0.000$ & $+0.015$ & \\
& NDCG@10 & $0.862$ & $0.878$ & $0.878 \pm 0.000$ & $+0.016$ & \\
\bottomrule
\end{tabular}
\end{table}

\emph{The protective layer is essentially free in the projected space.}
The difference (proposed) vs.\ (baseline\_proj) is within the $\pm 0.001$
5-seed CI for every metric and every encoder. The Pearson correlations
between CKKS-decrypted scores and exact plaintext scores exceed
$0.9999$ in all cases. The proxy criterion $\sigma_{rec} \geq 0.10$
holds uniformly.

\emph{The actual quality cost is paid by SVD truncation, not by the
protective wrapper.} The split
$\Delta_{\mathrm{SVD}} = \mathrm{baseline\_proj} - \mathrm{baseline\_dense}$
is $+0.004/+0.018/+0.014$ for retrieval-trained encoders with
$d \geq 768$ (e5-base / e5-large / bge-m3), $-0.080$ for the
compact e5-small ($d = 384$, half-truncation removes half the embedding
volume) and $-0.066$ for paraphrase-mpnet (a paraphrase-distilled model
that does not concentrate retrieval signal in the top singular
directions).

\emph{The end-to-end accuracy budget is met for retrieval-trained
$d \geq 768$.} The end-to-end criterion
$\Delta\mathrm{Acc@1} \geq -0.05$ relative to the raw baseline holds
for e5-base ($+0.003$), e5-large ($+0.018$) and bge-m3 ($+0.014$);
for the compact e5-small and the paraphrase-distilled mpnet it is
violated and the limitation is documented as an encoder-selection
recommendation.

\paragraph{Linear-denoiser side effect.} On the two
retrieval-trained encoders with $d = 1024$, half SVD truncation
\emph{improves} the four ranking metrics over the raw baseline
($+0.018$ Acc@1 on e5-large, $+0.014$ on bge-m3). This is consistent
with the interpretation of $V_k$-projection as a linear denoiser:
the trailing singular directions, dominated by fp16 inference noise,
get filtered out, and the ranking in $\mathrm{span}(V_k)$ is cleaner
than in the raw $\mathbb{R}^d$.

\begin{figure}[tb]
\centering
\includegraphics[width=\linewidth]{fig_exp5_per_seed.png}
\caption{Per-seed Acc@10 of the proposed pipeline relative to
baseline\_proj on five rotation seeds; the spread is below the 5-seed
CI of the baseline.}
\label{fig:exp5-per-seed}
\end{figure}

% ----------------------------------------------------------------
\subsection{Experiment 4: end-to-end client-server latency}\label{ssec:latency}
% ----------------------------------------------------------------

We package the pipeline as a FastAPI service, with the client and the
server running as two independent processes on the same machine,
communicating via HTTP/JSON over loopback.
The setup measures the realistic HTTP framing,
base64 serialisation of CKKS ciphertexts, FastAPI/uvicorn stack and JSON
(de-)serialisation, but without physical network latency.
The encoder is multilingual-e5-small ($d = 384$, $k = 192$,
$M = 48$, $K_{\mathrm{cands}} = 40$, $N_{\mathrm{docs}} = 10^6$),
averaged over $500$ queries $\times$ $5$ rotation seeds.

\paragraph{On-wire payload sizes.} For the same configuration the
measured HTTP payloads are: request body
$\approx 0.29$~MB (one CKKS query ciphertext, base64-encoded JSON,
inflated $\approx 1.37\times$ over the $0.21$~MB raw ciphertext);
response body $\approx 11.4$~MB ($K_{\mathrm{cands}} = 40$ score
ciphertexts, base64 JSON). One-off on-boarding transfers the public PQ
artefact ($48$~MB at $k=192$) and the CKKS public/Galois keys
($\approx 8$~MB); these are not part of the per-query budget. Reporting
sizes alongside $p_{95}$ is necessary because the dominant scaling
pressure for larger $K_{\mathrm{cands}}$ or distributed deployments is
the response payload, not CPU time.

\begin{table}[tb]
\centering
\caption{Latency decomposition at the canonical operating point
($k = 192$) on a $10^6$-document corpus. The $k$-dependent server-side
ct-pt rerank is measured on the $5$-seed integral run; the $k$-independent
client-side and HTTP stages are measured on the client-server
HTTP-loopback PoC. The encoder forward pass is independent of the defence
layer and excluded from the total.}
\label{tab:latency}
\small
\begin{tabular}{@{}lrl@{}}
\toprule
Stage & $p_{95}$ (ms) & Notes \\ \midrule
Encoder forward pass (client) & $\sim 200$ & encoder-dependent; excluded from total \\
Projection to $\mathrm{span}(V_k)$ (client) & $0.1$ & negligible \\
Local PQ search (client, faiss) & $22.8$ & $M = 48$, $K_{\mathrm{cands}} = 40$, $10^6$ docs \\
CKKS encryption (client) & $4.9$ & one rotated query, $N_{\mathrm{poly}} = 8192$ \\
\textbf{Server CKKS rerank, ct-pt} & $\mathbf{283}$ & $40$ ct-pt ops at $k = 192$ (5-seed integral run) \\
HTTP $+$ serialisation (req$+$resp) & $38$ & base64 ciphertexts over loopback \\
CKKS decrypt + sort (client) & $20.7$ & 40 score ciphertexts \\ \midrule
Total (excluding encoder) & $\approx 370$~ms & $p_{95}$, full pipeline \\
\bottomrule
\end{tabular}
\end{table}

\begin{figure}[tb]
\centering
\includegraphics[width=0.49\linewidth]{fig_exp6_latency_decomposition.png}\hfill
\includegraphics[width=0.49\linewidth]{fig_exp6_latency_cdf.png}
\caption{Left: per-stage latency decomposition at $k = 192$ (the
server-side ct-pt rerank dominates). Right: CDF of the server-side rerank
latency over $500$~queries $\times$ $5$~seeds ($p_{95} \approx 283$~ms);
the full end-to-end $p_{95} \approx 370$~ms (Table~\ref{tab:latency})
stays well below the $1$~s SLA.}
\label{fig:exp6-latency}
\end{figure}

% =====================================================================
\section{Security Analysis Under Known-Plaintext and Stronger Attackers}
\label{sec:stronger}
% =====================================================================

The experiments above establish the construction's behaviour \emph{inside}
its non-adaptive threat model (Section~\ref{sec:threat}). A privacy paper
that stopped there would be open to the objection that the only tested
adversary is also the weakest one. We therefore deliberately step outside
the model and measure four stronger channels---known-plaintext recovery of
the rotation, a text-level reference-corpus lookup, leakage from the public
PQ artefact, and an aligned off-the-shelf inverter---followed by a baseline
suite and a benchmark-generality check. Several of these are
\emph{negative} results for broad document-privacy claims; reporting them is
what licenses the deliberately narrow framing of
Sections~\ref{sec:discussion}--\ref{sec:conclusion}.

% ----------------------------------------------------------------
\subsection{Known-plaintext recovery of the secret rotation}
\label{ssec:procrustes}
% ----------------------------------------------------------------

The secret rotation $R$ is the only document-side layer beyond SVD
truncation. We test it against a known-plaintext attacker who holds $m$
pairs $(E(x_i), E_{\mathrm{rot},i})$ and estimates the orientation with the
orthogonal Procrustes solution~\cite{ref_schonemann}. The setup uses
\texttt{e5-small} projected to $k=192$, five rotation seeds, $400$ held-out
probes and a $10\,000$-document gallery; the attack is scored in embedding
space, where a successful alignment lets a probe recover its own projected
vector and retrieve the matching gallery document. This is the relevant
precursor to a few-shot alignment attack because it hands the inverter the
native projected space.

\begin{table}[tbp]
\centering
\caption{Known-plaintext Procrustes attack against the secret rotation
($m$ leaked pairs, means over five seeds). With $\approx k$ pairs the
orientation is numerically recovered.}
\label{tab:procrustes}
\small
\begin{tabular}{@{}rrrrr@{}}
\toprule
Known pairs $m$ & Mean cosine & Rel.\ $L_2$ error & Target R@1 & Target R@10 \\ \midrule
$0$   & $-0.004$ & $1.417$ & $0.000$ & $0.001$ \\
$10$  & $0.098$  & $1.343$ & $0.005$ & $0.039$ \\
$50$  & $0.377$  & $1.116$ & $0.722$ & $0.925$ \\
$100$ & $0.621$  & $0.870$ & $1.000$ & $1.000$ \\
$192$ & $1.000$  & $1.8{\times}10^{-5}$ & $1.000$ & $1.000$ \\
$500$ & $1.000$  & $4.9{\times}10^{-7}$ & $1.000$ & $1.000$ \\
\bottomrule
\end{tabular}
\end{table}

\begin{figure}[tbp]
\centering
\includegraphics[width=0.6\linewidth]{fig_exp7_procrustes.png}
\caption{Known-plaintext alignment destroys the unknown-rotation
assumption: with about $k$ leaked pairs the rotation is recovered. The
$x$ axis uses a symlog scale to include $m=0$.}
\label{fig:procrustes}
\end{figure}

The result is unambiguous (Table~\ref{tab:procrustes},
Fig.~\ref{fig:procrustes}): the rotation is useful only while aligned
plaintext anchors are unavailable. Even $50$ pairs---well below the
retained dimension $k=192$---already retrieve the target in the top~10 for
$92.5\%$ of probes in a $10\,000$-document gallery, and at $m\approx k$ the
orientation is numerically exact. The rotation therefore raises the bar
against an attacker with \emph{no} aligned plaintext, but it is not a
defence against known plaintext; this is reflected directly in the
threat-coverage table.

% ----------------------------------------------------------------
\subsection{Reference-corpus lookup after alignment}
\label{ssec:reference}
% ----------------------------------------------------------------

The Procrustes test is in embedding space. We next ask whether the same
failure becomes a \emph{text-level} channel without training any generative
inverter. The attacker is given the protected vectors, $m$ known pairs for
estimating $R$, and a reference corpus of candidate paragraphs with native
\texttt{e5-small} embeddings---modelling the realistic case where some
private documents also appear in a public or previously leaked corpus. We
use $500$ targets, $100\,000$ decoys, five seeds and the same $k=192$
subspace; the overlapping reference contains the targets, the disjoint
reference removes them and serves only as a lexical nearest-neighbour
proxy.

\begin{table}[tbp]
\centering
\caption{Reference-corpus lookup after known-plaintext alignment
(overlapping reference: $500$ targets $+$ $100\,000$ decoys; Jaccard is
token-set overlap with the top-1 retrieved paragraph; means over five
seeds).}
\label{tab:reference-attack}
\small
\begin{tabular}{@{}rrrrr@{}}
\toprule
Known pairs $m$ & Exact R@1 & Exact R@10 & Jaccard overlap-ref & Jaccard disjoint-ref \\ \midrule
$0$   & $0.000$ & $0.000$ & $0.006$ & $0.006$ \\
$10$  & $0.000$ & $0.009$ & $0.009$ & $0.009$ \\
$25$  & $0.044$ & $0.132$ & $0.055$ & $0.012$ \\
$50$  & $0.535$ & $0.784$ & $0.543$ & $0.018$ \\
$100$ & $0.998$ & $1.000$ & $0.999$ & $0.034$ \\
$192$ & $1.000$ & $1.000$ & $1.000$ & $0.043$ \\
\bottomrule
\end{tabular}
\end{table}

\begin{figure}[tbp]
\centering
\includegraphics[width=0.8\linewidth]{fig_exp9_reference_attack.png}
\caption{Reference-corpus lookup after alignment. When the reference
overlaps the protected collection, alignment turns the protected vector
into an exact paragraph lookup; with a disjoint reference this
token-overlap proxy stays low.}
\label{fig:reference-attack}
\end{figure}

The overlap case is severe (Table~\ref{tab:reference-attack}): with only
$50$ known pairs the attacker recovers the exact source paragraph at
$53.5\%$ top-1 among $100\,500$ candidates, rising to $99.8\%$ at $m=100$.
This is the text-level analogue of the Procrustes result and makes the
static-side limitation concrete: if a public or leaked reference corpus
overlaps the private collection, rotation secrecy is insufficient. The
disjoint-reference column is reported as a boundary condition---token
Jaccard stays at $0.034$ ($m=100$), so a non-overlapping reference needs
semantic metrics or a generative inverter before text reconstruction can
be claimed.

% ----------------------------------------------------------------
\subsection{Leakage from the public PQ artefact}
\label{ssec:pqleak}
% ----------------------------------------------------------------

Stage-1 candidate filtering publishes a product-quantization codebook plus
per-document codes, trained in the rotated space. We quantify how much of
$E_{\mathrm{rot}}$ this public artefact alone reconstructs on a
$20\,000$-document \texttt{e5-small} sample (Table~\ref{tab:pq-leakage}).
The canonical configuration ($M=48$, $8$ bits) preserves mean cosine
$0.953$ to $E_{\mathrm{rot}}$, $67.4\%$ of exact top-10 neighbours, and
$97.1\%$ of exact top-10 neighbours inside the approximate top-40
candidate set---that is, the public codes alone almost perfectly reproduce
the candidate set the server reranks.

\begin{table}[tbp]
\centering
\caption{Leakage from the public PQ codes (rotated \texttt{e5-small},
$20\,000$ documents; neighbour overlap excludes the query document).}
\label{tab:pq-leakage}
\small
\begin{tabular}{@{}rrrrr@{}}
\toprule
$M$ & Bits & Bytes/vector & Cosine to $E_{\mathrm{rot}}$ & NN overlap@10 \\ \midrule
$24$ & $8$ & $24$ & $0.837$ & $0.384$ \\
$48$ & $8$ & $48$ & $0.953$ & $0.674$ \\
$48$ & $6$ & $36$ & $0.904$ & $0.529$ \\
\bottomrule
\end{tabular}
\end{table}

These numbers do not weaken the CKKS query-privacy contribution, but they
do narrow the document-privacy interpretation: both the known-plaintext
anchor channel and the public PQ artefact must be treated as exposed. A
deployment that needs stronger static-side confidentiality requires either
a different static primitive or a composition that hides the PQ artefact
itself.

% ----------------------------------------------------------------
\subsection{Aligned off-the-shelf inversion stress test}
\label{ssec:aligned}
% ----------------------------------------------------------------

Finally we run a generative inverter against the aligned space. The
attacker is given GTR-base embeddings ($d=768$), SVD truncation to
$k=d/2$, the secret-rotation setting, and up to $m=500$ known pairs for
Procrustes alignment; the inverter is the off-the-shelf \textsc{Vec2Text}
corrector with a larger budget than Experiment~\ref{ssec:vec2text} ($24$
corrector iterations, max input length~$96$). For deterministic
reproducibility on a host that could not fetch external corpora the run
uses the script's offline synthetic-news+PII corpus ($120$ held-out
texts). The decisive control is the raw-embedding row.

\begin{table}[tbp]
\centering
\caption{Aligned off-the-shelf \textsc{Vec2Text} stress test ($120$
held-out texts; $m$ is the number of pairs used for alignment). Exact
match and typed-PII recall are zero in every row---but so is the raw
control, so this corrector cannot \emph{certify} protection.}
\label{tab:adaptive-inversion}
\small
\begin{tabular}{@{}lrrrr@{}}
\toprule
Case & Token F1 & BLEU & Exact & PII recall \\ \midrule
Raw embedding & $0.133$ & $0.0133$ & $0.000$ & $0.000$ \\
Known-$R$ oracle & $0.138$ & $0.0134$ & $0.000$ & $0.000$ \\
Unknown $R$ ($=$ Procrustes $m{=}0$) & $0.121$ & $0.0129$ & $0.000$ & $0.000$ \\
Procrustes, $m=50$ & $0.126$ & $0.0126$ & $0.000$ & $0.000$ \\
Procrustes, $m=100$ & $0.141$ & $0.0127$ & $0.000$ & $0.000$ \\
Procrustes, $m=192$ & $0.134$ & $0.0132$ & $0.000$ & $0.000$ \\
Procrustes, $m=500$ & $0.135$ & $0.0133$ & $0.000$ & $0.000$ \\
\bottomrule
\end{tabular}
\end{table}

The reading is honest but limited (Table~\ref{tab:adaptive-inversion}): no
exact document or typed PII is recovered even by a known-$R$ oracle, yet
the \emph{raw} embedding row is also at the BLEU floor ($0.013$), so this
off-the-shelf corrector is too weak to certify the transform. It does not
contradict the exact-lookup failure of Section~\ref{ssec:reference}; the
remaining open risk is a learned or corpus-adapted decoder trained
end-to-end on the rotated space, which we keep as the highest-priority
future experiment (Section~\ref{sec:limits}).

% ----------------------------------------------------------------
\subsection{Baseline suite: projection and calibrated noise}
\label{ssec:baselines}
% ----------------------------------------------------------------

To position the transform against natural alternatives rather than against
the raw system alone, we compare SVD projection with Gaussian random
projection and random orthogonal projection~\cite{ref_jl} at the same
$k=d/2$ (Table~\ref{tab:projection-baselines}), and with calibrated
isotropic Gaussian noise matched at equal mean relative distortion
(Table~\ref{tab:noise-sweep}).

\begin{table}[tbp]
\centering
\caption{Acc@1 under lightweight projection baselines at $k=d/2$ (random
baselines: mean $\pm$ s.d.\ over three seeds).}
\label{tab:projection-baselines}
\small
\begin{tabular}{@{}lrrrr@{}}
\toprule
Encoder & Raw & SVD & Gaussian RP & Orthogonal RP \\ \midrule
e5-small & $0.782$ & $0.702$ & $0.688 \pm 0.016$ & $0.737 \pm 0.010$ \\
e5-base  & $0.830$ & $0.834$ & $0.781 \pm 0.011$ & $0.804 \pm 0.017$ \\
mpnet    & $0.716$ & $0.650$ & $0.698 \pm 0.003$ & $0.705 \pm 0.008$ \\
e5-large & $0.796$ & $0.814$ & $0.765 \pm 0.010$ & $0.769 \pm 0.009$ \\
bge-m3   & $0.818$ & $0.832$ & $0.789 \pm 0.004$ & $0.810 \pm 0.002$ \\
\bottomrule
\end{tabular}
\end{table}

\begin{table}[tbp]
\centering
\caption{SVD truncation vs.\ independent Gaussian noise at matched mean
relative distortion ($100\,000$-document gallery). NN overlap@10 is the
overlap with raw-space neighbours; lower means less raw geometry
preserved.}
\label{tab:noise-sweep}
\scriptsize
\setlength{\tabcolsep}{4pt}
\begin{tabular}{@{}lrrrrrr@{}}
\toprule
Encoder & $k/d$ & $\sigma_{rec}$ & SVD Acc@1 & SVD NN@10 & Noise Acc@1 & Noise NN@10 \\ \midrule
e5-small & $0.125$ & $0.385$ & $0.318$ & $0.232$ & $0.750$ & $0.283$ \\
e5-small & $0.250$ & $0.332$ & $0.600$ & $0.333$ & $0.788$ & $0.367$ \\
e5-small & $0.500$ & $0.237$ & $0.814$ & $0.468$ & $0.828$ & $0.531$ \\
e5-small & $0.875$ & $0.079$ & $0.880$ & $0.551$ & $0.872$ & $0.830$ \\
e5-base  & $0.125$ & $0.354$ & $0.708$ & $0.370$ & $0.864$ & $0.503$ \\
e5-base  & $0.250$ & $0.268$ & $0.872$ & $0.501$ & $0.886$ & $0.641$ \\
e5-base  & $0.500$ & $0.129$ & $0.916$ & $0.578$ & $0.902$ & $0.829$ \\
e5-base  & $0.875$ & $0.023$ & $0.920$ & $0.592$ & $0.910$ & $0.969$ \\
\bottomrule
\end{tabular}
\end{table}

\begin{figure}[tbp]
\centering
\includegraphics[width=0.8\linewidth]{fig_exp8_tradeoff.png}
\caption{Utility/leakage diagnostic for SVD truncation and calibrated
Gaussian noise (dashed: matched-$\sigma_{rec}$ noise baselines).}
\label{fig:noise-sweep}
\end{figure}

Two engineering conclusions follow. First, the projection choice is
encoder-dependent: SVD is strongest on the retrieval-trained $d\geq768$
encoders that are the method's operating region, but random orthogonal
projection beats SVD on \texttt{mpnet} and \texttt{e5-small}---so a
deployment should run a baseline sweep, not assume SVD is universally best.
Second, $\sigma_{rec}$ is not a privacy metric: at $k/d=0.875$ matched
noise preserves almost all raw neighbours, and at the main $k/d=0.5$
operating point SVD removes far more raw-neighbour structure than matched
noise ($0.578$ vs.\ $0.829$ NN overlap@10 on \texttt{e5-base}) while
keeping Acc@1. The defensible claim is thus the narrow one: SVD is an
encoder-dependent utility/leakage trade-off, not a noise mechanism that
dominates all alternatives.

% ----------------------------------------------------------------
\subsection{Benchmark generality on BEIR}
\label{ssec:beir}
% ----------------------------------------------------------------

The integral experiment uses one-million-document self-retrieval, a
controlled geometry probe. To check that the metric-preservation reading
is not an artefact of self-retrieval, we repeat the raw-vs-SVD comparison
on two standard BEIR~\cite{ref_beir} datasets with graded relevance
judgements---SciFact ($5{,}183$ docs, $300$ queries) and NFCorpus ($3{,}633$
docs, $323$ queries)---fitting $V_k$ on the corpus at $k=d/2$ and comparing
exact inner-product retrieval in $\mathbb{R}^d$ and in
$\mathrm{span}(V_k)$.

\begin{table}[tbp]
\centering
\caption{SVD truncation ($k=d/2$) vs.\ the raw space on BEIR.
$\Delta=\mathrm{SVD}-\mathrm{raw}$; brackets are paired-bootstrap $95\%$
CIs ($10^4$ resamples). \textbf{Every} $\Delta$ nDCG@10 CI excludes~$0$.}
\label{tab:beir}
\footnotesize
\setlength{\tabcolsep}{4pt}
\begin{tabular}{@{}llrrrr@{}}
\toprule
Encoder & Dataset & $\sigma_{rec}$ & raw nDCG@10 & SVD nDCG@10 & $\Delta$ nDCG@10 (95\% CI) \\
\midrule
e5-large & SciFact  & $0.072$ & $0.643$ & $0.574$ & $-0.069\,[-0.091,-0.047]$ \\
e5-base  & SciFact  & $0.090$ & $0.637$ & $0.586$ & $-0.051\,[-0.072,-0.032]$ \\
e5-small & SciFact  & $0.183$ & $0.598$ & $0.519$ & $-0.079\,[-0.103,-0.055]$ \\
e5-base  & NFCorpus & $0.085$ & $0.327$ & $0.305$ & $-0.022\,[-0.032,-0.013]$ \\
e5-small & NFCorpus & $0.172$ & $0.302$ & $0.277$ & $-0.025\,[-0.035,-0.014]$ \\
\bottomrule
\end{tabular}
\end{table}

Table~\ref{tab:beir} is a deliberate corrective to the self-retrieval
reading. On both graded-relevance datasets and for every encoder, half-SVD
truncation \emph{significantly reduces} nDCG@10 (all CIs exclude zero,
drops of $2$--$8$ points). Decisively, \texttt{e5-large}---the one encoder
that \emph{gained} on the $10^6$ self-retrieval probe---\emph{loses}
$0.069$ nDCG@10 on SciFact. The denoiser effect is therefore real but
\emph{task-dependent}: on real graded IR the discarded spectrum still
carries relevance, so $k$ must be treated as a tunable accuracy cost via a
$k$-sweep, exactly as the encoder-selection discussion below recommends.
(Full end-to-end CKKS reranking on BEIR is left to future work; this
experiment isolates the projection-fidelity property that the reranking
preserves.)

% =====================================================================
\section{Discussion}\label{sec:discussion}
% =====================================================================

\paragraph{Encoder selection is part of the protocol.} The protective
wrapper is essentially metric-preserving in the projected space; the
end-to-end deviation from the raw baseline is dominated by what
SVD truncation does on the chosen encoder. Models trained with
contrastive retrieval objectives (\textsc{e5-*}, \textsc{bge-*})
concentrate the retrieval signal in the top singular directions and
recover most of the original quality after $k = d/2$ truncation; some of
them even \emph{benefit} from the truncation (linear-denoiser effect).
Paraphrase-distilled models (\textsc{mpnet}) and very compact models
(\textsc{e5-small}) do not show this concentration, and we recommend
either replacing them with a retrieval-trained equivalent or running a
$k$-sweep before deployment to find a $k$ that simultaneously satisfies
the application's accuracy budget and $\sigma_{rec} \geq 0.10$.

\paragraph{The role of the rotation $R$: obfuscation, not a primitive.}
The secret rotation is an empirical obfuscation layer, not a
cryptographic primitive. It contributes a defence channel orthogonal to
SVD truncation, and even at $k/d = 1.0$ the unknown-$R$ regime reduces
off-the-shelf Vec2Text BLEU to the noise floor. Its limits are now
measured rather than asserted (Section~\ref{sec:stronger}): a known-$R$
attacker gains nothing beyond SVD; a known-plaintext attacker recovers
$R$ by orthogonal Procrustes from $\approx k$ leaked pairs
(Table~\ref{tab:procrustes}), which becomes an exact paragraph lookup
when a reference corpus overlaps the collection
(Table~\ref{tab:reference-attack}); and an end-to-end learned decoder
remains the main untested case. The rotation's value is therefore that
it is the cheapest architectural layer with a measurable effect on an
attacker who lacks aligned plaintext---not a security guarantee. The
contribution of this paper is the measured privacy/latency/accuracy
trade-off and its quantified boundaries, not the rotation in isolation.

\paragraph{What CKKS hides and what it does not.} The CKKS layer hides
the numerical representation of the query and the exact similarity
scores from the server. Two things it does \emph{not} hide are the
identifiers of the candidates produced by stage-1 PQ filtering and the
$L_2$-distances among them implicit in the order in which they are
sent. An access-pattern attacker can therefore, over many queries,
build a co-occurrence model of which documents are co-retrieved.
Composing the pipeline with a PIR-style retrieval primitive (Tiptoe
\cite{ref_tiptoe}, SealPIR \cite{ref_sealpir}, OnionPIR
\cite{ref_onionpir}, SimplePIR \cite{ref_simplepir}) is the natural
remedy; ours is engineered to be PIR-friendly because the
ct-pt operations on the server are stateless and cluster-local.

\paragraph{Parameter selection vs.\ hand-tuning.} Compared with hand-tuned
TenSEAL stock parameters, the configuration chosen by the offline grid search
(Section~\ref{sec:autotune}) achieves a $\approx 1.7\times$ speed-up at the
same security level and the same ranking quality (within
$\Delta \mathrm{Acc@1} \leq 1$~p.p.\ of the SVD-projected baseline). The
selection is a one-off offline cost and adds nothing at run time; porting to a
different hardware platform only requires re-running the same small benchmark
grid.

% =====================================================================
\section{Limitations and Required Experiments}\label{sec:limits}
% =====================================================================

We separate \emph{intrinsic limitations} (properties of the construction
under study) from the \emph{remaining open experiments}. The
stronger-attacker, baseline and generality experiments that earlier drafts
listed as mandatory future work are now reported in
Section~\ref{sec:stronger}; what is listed below as open is the genuinely
residual set, kept as named scoped experiments rather than vague
``future work'' as a deliberate hedge against overclaiming.

\paragraph{Intrinsic limitations.}
\begin{itemize}[leftmargin=18pt,topsep=2pt,itemsep=2pt]
\item \emph{Document privacy is not cryptographic.} $E_{\mathrm{rot}}$
is plaintext on the server; it is protected only by SVD truncation (a
proxy) and the secret rotation (an empirical obfuscation layer). A
known-$R$, known-plaintext or adaptive adversary defeats this layer.
\item \emph{Access patterns and PQ codes are exposed.} CKKS hides values,
not which candidates are reranked, and the public PQ codes are a lossy
view of $E_{\mathrm{rot}}$. Production use requires composition with a
PIR/ORAM layer plus a PQ-leakage analysis.
\item \emph{Encoder-dependence of the accuracy budget.} The strict
end-to-end criterion $\Delta\mathrm{Acc@1} \geq -5$~p.p.\ holds only
on retrieval-trained encoders with $d \geq 768$. Compact or
paraphrase-distilled encoders need either replacement or a less
aggressive $k$ (which then violates the $\sigma_{rec} \geq 0.10$ proxy).
\item \emph{Single-machine evaluation.} The latency numbers in
Experiment~\ref{ssec:latency} are taken on loopback HTTP; a
geographically distributed deployment will add (uncharacterised)
network latency.
\end{itemize}

\paragraph{Completed in this revision.} The stronger-attacker programme is
now executed (Section~\ref{sec:stronger}): known-plaintext Procrustes
recovery of $R$ (Table~\ref{tab:procrustes}); the text-level
reference-corpus lookup (Table~\ref{tab:reference-attack}); public-PQ-code
leakage (Table~\ref{tab:pq-leakage}); an aligned off-the-shelf inverter
(Table~\ref{tab:adaptive-inversion}); a projection and calibrated-noise
baseline suite covering random projection, SVD-without-rotation and
matched-distortion Gaussian noise
(Tables~\ref{tab:projection-baselines}--\ref{tab:noise-sweep}); and a BEIR
generality check on the projection step (Table~\ref{tab:beir}).

\paragraph{Remaining open experiments.}
\begin{enumerate}[leftmargin=20pt,topsep=2pt,itemsep=3pt]
\item \emph{End-to-end learned decoder (highest priority).} The aligned
off-the-shelf corrector of Section~\ref{ssec:aligned} is too weak to
\emph{certify} protection---its raw control is already at the BLEU floor.
The decisive test is an inversion model \emph{trained} on rotated/SVD
embeddings of part of the corpus and evaluated on a holdout; the current
narrative explicitly does not claim robustness against it.
\item \emph{Typed-PII, multi-encoder $\sigma_{rec}$ ablation.} Replace the
single-encoder, BLEU-only calibration of the $\sigma_{rec}=0.10$ proxy with
a multi-encoder ablation reporting BLEU, token overlap and typed-PII
recovery (names, addresses, e-mail, phone, medical terms), which is what
justifies (or refutes) the threshold.
\item \emph{Modern attack suite and full-pipeline benchmarks.} Repeat the
inversion study against the 2024--2026 attack generations that exploit the
universal geometry of embedding spaces rather than a fixed basis
orientation---few-shot alignment (\textsc{ALGEN}~\cite{ref_algen}),
zero-shot inversion (\textsc{ZSInvert}~\cite{ref_zsinvert}), unsupervised
cross-space translation (\textsc{vec2vec}~\cite{ref_vec2vec}) and online
cross-domain inversion (\textsc{Zero2Text}~\cite{ref_zero2text})---with real
text, and run the \emph{complete} CKKS reranking (not only the projection
step measured in Table~\ref{tab:beir}) on BEIR, MIRACL and MS~MARCO.
\end{enumerate}

% =====================================================================
\section{Conclusion}\label{sec:conclusion}
% =====================================================================

This paper is a scientific study of how SVD truncation and an unknown
secret rotation reshape the privacy--accuracy--latency trade-off of
semantic search; it is not a systems-engineering report and does not
claim a deployable system. Three findings answer the research questions.
\textbf{(RQ1)} A formal lower-$L_2$ bound (Lemma~\ref{lemma:proj-l2})
characterises exactly what a projection-restricted decoder cannot
recover; it is a projection bound, not an inversion-security theorem,
and it makes the \emph{proxy} criterion $\sigma_{rec} \geq 0.10$ precise
only within that decoder class, the threshold itself being an
engineering proxy. \textbf{(RQ2)} At $10^6$-document scale the
construction's protective layer is metric-preserving in
$\mathrm{span}(V_k)$ (within the 5-seed CI), the query side is
cryptographically protected by CKKS at sub-second $p_{95}$, while
document and access-pattern privacy are \emph{not} cryptographic and
the public PQ artefact is an exposed channel whose leakage we now
quantify (Section~\ref{sec:stronger})---so the secret rotation is
characterised as an empirical obfuscation effect against an off-the-shelf
inverter, not as a guarantee. \textbf{(RQ3)} For
retrieval-trained encoders with $d \geq 768$, data-driven SVD truncation
to $k = d/2$ acts as a linear \emph{denoiser}, improving the four
ranking metrics over the raw baseline rather than degrading them---a
finding of independent interest beyond the privacy setting.

We do not claim the privacy problem is solved. This work delivers a
measured privacy, accuracy and latency operating point and a characterised
denoising phenomenon under a clearly delimited threat model, together
with a stronger-attacker analysis (Section~\ref{sec:stronger}) that
quantifies where the document-side obfuscation fails---known-plaintext
recovery of $R$, reference-corpus lookup and public-PQ-code leakage---and
leaves an end-to-end learned decoder as the principal open question.
Composing the protocol with a PIR-style access-pattern hider---natural
given the stateless server-side ct-pt operations---is the most attractive
architectural extension.

% =====================================================================
\section*{Reproducibility}
% =====================================================================

All scripts, notebooks, configuration files and JSON output of the
experiments are released at
\url{https://github.com/sergkurilenko/research/tree/article1}.
The integral experiment is driven by a single Python script with
deterministic seeds and a JSON output file; the CKKS parameter-selection
benchmark ships its configuration grid and timings as a CSV. The corpus
(one-million Russian-Wikipedia paragraphs) is reproducibly sliced from a
public Wikipedia dump with the date and content checksum recorded.

% =====================================================================
\begin{thebibliography}{99}\small
% =====================================================================

\bibitem{ref_sbert} N.\ Reimers and I.\ Gurevych, ``Sentence-BERT:
Sentence Embeddings using Siamese BERT-Networks,'' in
\emph{Proc.\ EMNLP}, 2019, pp.\ 3982--3992.

\bibitem{ref_e5} L.\ Wang, N.\ Yang, X.\ Huang, et al.,
``Text Embeddings by Weakly-Supervised Contrastive Pre-training,''
\emph{arXiv:2212.03533}, 2022.

\bibitem{ref_gtr} J.\ Ni et al., ``Large Dual Encoders Are Generalizable
Retrievers,'' in \emph{Proc.\ EMNLP}, 2022, pp.\ 9844--9855.

\bibitem{ref_bgem3} J.\ Chen, S.\ Xiao, P.\ Zhang, K.\ Luo, D.\ Lian,
and Z.\ Liu, ``BGE M3-Embedding: Multi-Lingual, Multi-Functionality,
Multi-Granularity Text Embeddings Through Self-Knowledge Distillation,''
\emph{arXiv:2402.03216}, 2024.

\bibitem{ref_dpr} V.\ Karpukhin et al., ``Dense Passage Retrieval for
Open-Domain Question Answering,'' in \emph{Proc.\ EMNLP}, 2020,
pp.\ 6769--6781.

\bibitem{ref_colbert} O.\ Khattab and M.\ Zaharia, ``ColBERT: Efficient
and Effective Passage Search via Contextualised Late Interaction over
BERT,'' in \emph{Proc.\ ACM SIGIR}, 2020, pp.\ 39--48.

\bibitem{ref_vec2text} J.\ X.\ Morris, V.\ Kuleshov, V.\ Shmatikov, and
A.\ M.\ Rush, ``Text Embeddings Reveal (Almost) As Much As Text,'' in
\emph{Proc.\ EMNLP}, 2023.

\bibitem{ref_geia} H.\ Li, M.\ Xu, and Y.\ Song, ``Sentence Embedding
Leaks More Information than You Expect: Generative Embedding Inversion
Attack to Recover the Whole Sentence,'' in \emph{Findings of ACL}, 2023,
pp.\ 14022--14040.

\bibitem{ref_teia} Y.-H.\ Huang, Y.\ Tsai, H.\ Hsiao, H.-Y.\ Lin, and
S.-D.\ Lin, ``Transferable Embedding Inversion Attack: Uncovering Privacy
Risks in Text Embeddings without Model Queries,'' in \emph{Proc.\ ACL
(Long)}, 2024, pp.\ 4193--4205.

\bibitem{ref_algen} Y.\ Chen, Q.\ Xu, and J.\ Bjerva, ``ALGEN: Few-shot
Inversion Attacks on Textual Embeddings via Cross-Model Alignment and
Generation,'' in \emph{Proc.\ ACL (Long)}, 2025, pp.\ 24330--24348.
\emph{arXiv:2502.11308}.

\bibitem{ref_zsinvert} C.\ Zhang, J.\ X.\ Morris, and V.\ Shmatikov,
``Universal Zero-shot Embedding Inversion,'' \emph{arXiv:2504.00147}, 2025.

\bibitem{ref_vec2vec} R.\ Jha, C.\ Zhang, V.\ Shmatikov, and J.\ X.\ Morris,
``Harnessing the Universal Geometry of Embeddings (vec2vec),''
\emph{arXiv:2505.12540}, 2025.

\bibitem{ref_zero2text} D.\ Kim, D.\ Kang, K.\ Lee, H.\ Baek, and
B.\ B.\ Kang, ``Zero2Text: Zero-Training Cross-Domain Inversion Attacks on
Textual Embeddings,'' \emph{arXiv:2602.01757}, 2026.

\bibitem{ref_zeng_rag} S.\ Zeng et al., ``The Good and the Bad:
Exploring Privacy Issues in Retrieval-Augmented Generation (RAG),''
in \emph{Findings of ACL}, 2024, pp.\ 4505--4524.

\bibitem{ref_song_leak} C.\ Song and A.\ Raghunathan, ``Information
Leakage in Embedding Models,'' in \emph{Proc.\ ACM CCS}, 2020,
pp.\ 377--390.

\bibitem{ref_shokri} R.\ Shokri et al., ``Membership Inference Attacks
Against Machine Learning Models,'' in \emph{Proc.\ IEEE S\&P},
2017, pp.\ 3--18.

\bibitem{ref_carlini} N.\ Carlini et al., ``Extracting Training Data
from Large Language Models,'' in \emph{Proc.\ USENIX Security},
2021, pp.\ 2633--2650.

\bibitem{ref_dwork} C.\ Dwork, ``Differential Privacy,'' in
\emph{Proc.\ ICALP}, 2006, pp.\ 1--12.

\bibitem{ref_dpsgd} M.\ Abadi et al., ``Deep Learning with Differential
Privacy,'' in \emph{Proc.\ ACM CCS}, 2016, pp.\ 308--318.

\bibitem{ref_lyu} L.\ Lyu, X.\ He, and Y.\ Li, ``Differentially Private
Representation for NLP,'' in \emph{Findings of EMNLP}, 2020,
pp.\ 2355--2365.

\bibitem{ref_kenthapadi} K.\ Kenthapadi, A.\ Korolova, I.\ Mironov, and
N.\ Mishra, ``Privacy via the Johnson-Lindenstrauss Transform,''
\emph{J.\ Privacy and Confidentiality}, vol.\ 5, no.\ 1, 2013.

\bibitem{ref_liu} K.\ Liu, H.\ Kargupta, and J.\ Ryan,
``Random Projection-Based Multiplicative Data Perturbation for Privacy
Preserving Distributed Data Mining,'' \emph{IEEE TKDE}, vol.\ 18,
no.\ 1, 2006, pp.\ 92--106.

\bibitem{ref_ckks} J.\ H.\ Cheon, A.\ Kim, M.\ Kim, and Y.\ Song,
``Homomorphic Encryption for Arithmetic of Approximate Numbers,'' in
\emph{Advances in Cryptology---ASIACRYPT 2017}, LNCS 10624, pp.\ 409--437.

\bibitem{ref_ckks_bootstrap} J.\ H.\ Cheon, K.\ Han, A.\ Kim et al.,
``Bootstrapping for Approximate Homomorphic Encryption,'' in
\emph{Advances in Cryptology---EUROCRYPT 2018}, LNCS 10820,
pp.\ 360--384.

\bibitem{ref_he_std} M.\ Albrecht et al., ``Homomorphic Encryption
Security Standard,'' HomomorphicEncryption.org, 2018.

\bibitem{ref_lwe_estimator} M.\ R.\ Albrecht, R.\ Player, and S.\ Scott,
``On the Concrete Hardness of Learning with Errors,'' \emph{Journal of
Mathematical Cryptology}, vol.\ 9, no.\ 3, 2015, pp.\ 169--203
(lattice-estimator methodology).

\bibitem{ref_schonemann} P.\ H.\ Sch\"onemann, ``A Generalized Solution
of the Orthogonal Procrustes Problem,'' \emph{Psychometrika}, vol.\ 31,
no.\ 1, 1966, pp.\ 1--10.

\bibitem{ref_openfhe} A.\ Al Badawi et al., ``OpenFHE: Open-Source Fully
Homomorphic Encryption Library,'' in \emph{Proc.\ WAHC '22},
2022, pp.\ 53--63.

\bibitem{ref_eva} R.\ Dathathri et al., ``EVA: An Encrypted Vector
Arithmetic Language and Compiler for Efficient Homomorphic
Computation,'' in \emph{Proc.\ ACM PLDI}, 2020, pp.\ 546--561.

\bibitem{ref_chet} R.\ Dathathri et al., ``CHET: An Optimizing Compiler
for Fully-Homomorphic Neural-Network Inferencing,'' in
\emph{Proc.\ ACM PLDI}, 2019, pp.\ 142--156.

\bibitem{ref_jung_gpu} W.\ Jung, S.\ Kim, J.\ H.\ Ahn et al.,
``Over 100x Faster Bootstrapping in Fully Homomorphic Encryption
through Memory-centric Optimisation with GPUs,'' \emph{IACR TCHES},
vol.\ 2021, no.\ 4, pp.\ 114--148.

\bibitem{ref_hexl} F.\ Boemer et al., ``Intel HEXL: Accelerating
Homomorphic Encryption with Intel AVX512-IFMA52,'' in
\emph{Proc.\ WAHC '21}, 2021, pp.\ 57--62.

\bibitem{ref_cryptonets} R.\ Gilad-Bachrach et al., ``CryptoNets:
Applying Neural Networks to Encrypted Data with High Throughput and
Accuracy,'' in \emph{Proc.\ ICML}, 2016, vol.\ 48, pp.\ 201--210.

\bibitem{ref_tiptoe} A.\ Henzinger, E.\ Dauterman, H.\ Corrigan-Gibbs,
and N.\ Zeldovich, ``Private Web Search with Tiptoe,'' in
\emph{Proc.\ ACM SOSP}, 2023.

\bibitem{ref_sealpir} S.\ Angel, H.\ Chen, K.\ Laine, and S.\ Setty,
``PIR with Compressed Queries and Amortised Query Processing,'' in
\emph{Proc.\ IEEE S\&P}, 2018, pp.\ 962--979.

\bibitem{ref_simplepir} A.\ Henzinger, M.\ M.\ Hong, H.\ Corrigan-Gibbs,
S.\ Meiklejohn, and V.\ Vaikuntanathan, ``One Server for the Price of
Two: Simple and Fast Single-Server Private Information Retrieval
(SimplePIR),'' in \emph{Proc.\ USENIX Security}, 2023.

\bibitem{ref_onionpir} M.\ H.\ Mughees, H.\ Chen, and L.\ Ren,
``OnionPIR: Response Efficient Single-Server PIR,'' in
\emph{Proc.\ ACM CCS}, 2021, pp.\ 2292--2306.

\bibitem{ref_halko} N.\ Halko, P.-G.\ Martinsson, and J.\ A.\ Tropp,
``Finding Structure with Randomness: Probabilistic Algorithms for
Constructing Approximate Matrix Decompositions,'' \emph{SIAM Review},
vol.\ 53, no.\ 2, 2011, pp.\ 217--288.

\bibitem{ref_pq} H.\ Jegou, M.\ Douze, and C.\ Schmid, ``Product
Quantisation for Nearest Neighbor Search,'' \emph{IEEE TPAMI},
vol.\ 33, no.\ 1, 2011, pp.\ 117--128.

\bibitem{ref_hnsw} Y.\ A.\ Malkov and D.\ A.\ Yashunin, ``Efficient and
Robust Approximate Nearest Neighbour Search Using Hierarchical
Navigable Small World Graphs,'' \emph{IEEE TPAMI}, vol.\ 42,
no.\ 4, 2020, pp.\ 824--836.

\bibitem{ref_milvus} J.\ Wang, X.\ Yi, R.\ Guo et al., ``Milvus:
A Purpose-Built Vector Data Management System,'' in
\emph{Proc.\ ACM SIGMOD}, 2021, pp.\ 2614--2627.

\bibitem{ref_eckart} C.\ Eckart and G.\ Young, ``The Approximation of
One Matrix by Another of Lower Rank,'' \emph{Psychometrika}, vol.\ 1,
no.\ 3, 1936, pp.\ 211--218.

\bibitem{ref_jl} W.\ B.\ Johnson and J.\ Lindenstrauss, ``Extensions
of Lipschitz Mappings into a Hilbert Space,'' \emph{Contemporary
Mathematics}, vol.\ 26, 1984, pp.\ 189--206.

\bibitem{ref_beir} N.\ Thakur, N.\ Reimers, A.\ R\"uckl\'e, A.\ Srivastava,
and I.\ Gurevych, ``BEIR: A Heterogeneous Benchmark for Zero-shot
Evaluation of Information Retrieval Models,'' in \emph{Proc.\ NeurIPS
Datasets and Benchmarks Track}, 2021.

\bibitem{ref_bleu} K.\ Papineni et al., ``BLEU: a Method for Automatic
Evaluation of Machine Translation,'' in \emph{Proc.\ ACL}, 2002,
pp.\ 311--318.

\bibitem{ref_rag} P.\ Lewis et al., ``Retrieval-Augmented Generation
for Knowledge-Intensive NLP Tasks,'' in \emph{Proc.\ NeurIPS}, 2020.

\bibitem{ref_self_rag} A.\ Asai et al., ``Self-RAG: Learning to
Retrieve, Generate, and Critique through Self-Reflection,'' in
\emph{Proc.\ ICLR}, 2024.

\bibitem{ref_crag} S.-Q.\ Yan et al., ``Corrective Retrieval Augmented
Generation,'' \emph{arXiv:2401.15884}, 2024.

\bibitem{ref_rag_survey} Y.\ Gao et al., ``Retrieval-Augmented
Generation for Large Language Models: A Survey,'' \emph{arXiv:2312.10997},
2024.

\end{thebibliography}

\end{document}
